\section{OBJETIVOS}

\subsection{Objetivo General}

Desarrollar un metabuscador semántico de alojamientos temporales inspirado en
plataformas de intermediación tipo Airbnb, basado en una ontología OWL del
dominio, integrado con la base de conocimiento enlazada DBpedia e incorporando
soporte multilingüe mediante anotaciones \texttt{rdfs:label}, con el fin de
demostrar la capacidad de las tecnologías de la Web Semántica para la
representación y recuperación expresiva de información en dominios complejos.

\subsection{Objetivos Específicos}

\begin{itemize}

      \item Modelar formalmente el dominio de los alojamientos temporales mediante una
      ontología OWL que represente propiedades de hospedaje, amenidades, zonas
      geográficas, perfiles de huésped, propósitos de viaje y políticas de reserva,
      estableciendo relaciones semánticas y restricciones lógicas entre los
      conceptos del sistema.

      \item Diseñar un conjunto de preguntas de competencia que representen los
      principales escenarios de búsqueda dentro del dominio y que permitan evaluar
      la capacidad de la ontología para responder consultas complejas mediante el
      lenguaje SPARQL.

      \item Implementar consultas SPARQL capaces de resolver las preguntas de
      competencia definidas, combinando múltiples criterios de búsqueda relacionados
      con ubicación geográfica, características del alojamiento, disponibilidad de
      amenidades y contexto del viaje.

      \item Integrar la ontología con la base de conocimiento enlazada DBpedia con el
      fin de incorporar información geográfica real sobre ciudades, zonas urbanas y
      puntos de interés del territorio boliviano, enriqueciendo así la base de
      conocimiento del sistema.

      \item Incorporar soporte multilingüe dentro de la ontología mediante
      anotaciones \texttt{rdfs:label} con etiquetas de idioma, permitiendo que los
      conceptos, propiedades e instancias del modelo puedan representarse y
      recuperarse en diferentes lenguajes.

      \item Verificar la consistencia lógica del modelo ontológico utilizando un
      razonador OWL integrado en Protégé, garantizando la coherencia de la jerarquía
      de clases, las relaciones semánticas y las restricciones definidas en el
      modelo.

      \item Implementar un prototipo de aplicación web que permita a los usuarios
      explorar y buscar alojamientos temporales mediante una interfaz similar a las
      plataformas de reservas en línea o marketplaces de hospedaje. La interfaz
      permitirá visualizar propiedades, aplicar filtros de búsqueda y consultar
      resultados de forma intuitiva, mientras que el procesamiento de las consultas
      se realizará internamente sobre la base de conocimiento semántica mediante
      consultas SPARQL.

\end{itemize}